\chapter*{General conclusion }
\markboth{General conclusion}{} %pour afficher l'entete
\addcontentsline{toc}{chapter}{General conclusion }

This project has addressed the design, implementation, and validation of a simultaneous multifrequency Electrical Impedance Tomography (EIT) acquisition system dedicated to biomedical applications. The developed platform, built around the STM32G474 microcontroller, successfully generated and acquired multifrequency excitation signals while ensuring reliable current and voltage measurements through carefully designed conditioning circuits. Signal processing methods such as Goertzel filtering and Fourier analysis were implemented to extract impedance values with high accuracy, and the system was validated through simulations, hardware experiments, and experimental measurements. The reproducibility tests confirmed that the acquisition chain provided stable results across repeated trials and different excitation frequencies, while image reconstruction experiments demonstrated the capability of the system to localize conductive inclusions inside a water test-tank. These results underline the feasibility of the proposed approach and highlight the advantages of simultaneous multifrequency operation compared to conventional sequential techniques.

Nevertheless, some limitations were encountered. At higher excitation frequencies, electrode polarization and parasitic effects introduced variability in the measurements, while the reconstructed images exhibited blurring and distortions that can be attributed to the intrinsic ill-posedness of the EIT inverse problem and to the simplified reconstruction algorithms employed. Despite these challenges, the overall system performance demonstrates the potential of simultaneous multifrequency EIT as a compact, low-cost, and versatile imaging modality for biomedical monitoring.

Looking ahead, future perspectives include the refinement of both hardware and software components in order to achieve higher accuracy and resolution. On the hardware side, improvements such as optimized analog front-end design, the use of higher-performance converters, and parallel multi-channel acquisition would allow better noise immunity and faster frame rates. On the algorithmic side, more advanced reconstruction methods, including nonlinear solvers and machine learning–based approaches, could enhance image quality and robustness against measurement uncertainties. Extending the frequency range of operation would further enable a more comprehensive characterization of tissue properties, while the integration of real-time processing pipelines would bring the system closer to clinical applicability. Finally, moving from controlled test-tank experiments toward in-vitro and in-vivo validation will be a crucial step in demonstrating the system’s practical relevance in medical diagnostics, rehabilitation monitoring, and human–machine interaction.

In conclusion, this work has established a functional proof of concept for a simultaneous multifrequency EIT system, demonstrating its reproducibility, robustness, and potential for biomedical imaging. By building upon the foundations laid in this thesis, future research can advance toward more accurate, reliable, and real-time EIT platforms capable of contributing meaningfully to next-generation healthcare technologies.